\chapter{CHIA SẺ: MENTORING TRONG KỶ NGUYÊN AI \& VIBE CODE}

\section{Lời tự sự của ``Mentor trẻ nhất'' sau 3 mùa đồng hành}

Bước vào mùa thứ 3 của IMP, khi nhìn lại hành trình, tôi nhận ra mình có một xuất phát điểm rất khác biệt so với các anh chị Mentor kỳ cựu:
\begin{itemize}
    \item \textbf{Về trải nghiệm}: Ít hơn các Mentor lớn tuổi hàng chục năm ``chinh chiến'' thương trường.
    \item \textbf{Về nguồn lực}: Doanh nghiệp nhỏ, ngân sách hạn hạn, nhân sự mỏng.
    \item \textbf{Về công cụ}: Không có bộ máy back-office hùng hậu để hỗ trợ.
\end{itemize}

Tuy nhiên, ``trong cái khó ló cái khôn''. Chính sự thiếu hụt về nguồn lực truyền thống đã thúc đẩy tôi tìm đến một nguồn lực vô tận mới: \textbf{Trí tuệ nhân tạo (AI)}. Bài chia sẻ này đúc kết cách một Mentor trẻ ``vượt rào'' kinh nghiệm bằng công nghệ.

\section{``Trending Vibe Code'': Khi tư duy quan trọng hơn cú pháp}

Khái niệm ``Vibe Coding'' không chỉ đơn thuần là dùng AI viết code, mà là một sự chuyển dịch tư duy:
\begin{itemize}
    \item \textbf{Từ ``How'' sang ``What''}: Mentor không còn dạy Mentee cách viết từng dòng code, mà dạy cách ``vibe'' (cảm nhận) kiến trúc hệ thống và đặt câu hỏi đúng cho AI.
    \item \textbf{Ứng dụng trong Mentoring}: Tôi khuyến khích Mentee sử dụng các công cụ như Cursor, Windsurf để hiện thực hóa ý tưởng ngay lập tức. Điều này giúp Mentee không bị nản lòng bởi các lỗi cú pháp nhỏ, mà tập trung vào logic kinh doanh và giá trị sản phẩm.
\end{itemize}

\section{Hệ sinh thái AI trong vận hành và phát triển}

Để giải bài toán ``ít ngân sách, ít người'', tôi đã tích cực ứng dụng AI vào 3 trụ cột cốt lõi mà tôi cũng đang hướng dẫn cho các Mentee:

\subsection{AI trong Marketing (Làm nhiều hơn với ít hơn)}
\begin{itemize}
    \item \textbf{Sáng tạo nội dung}: Dùng AI để lên kế hoạch content cho cả tháng chỉ trong 30 phút.
    \item \textbf{Thiết kế hình ảnh}: Sử dụng các công cụ tạo ảnh (Generative AI) để xây dựng bộ nhận diện thương hiệu chuyên nghiệp mà không cần thuê designer ngoài.
    \item \textbf{Tự động hóa}: Dùng AI để phân tích insight khách hàng từ dữ liệu thô.
\end{itemize}

\subsection{AI trong Quản trị dự án (Chuyên nghiệp hóa dòng chảy)}
\begin{itemize}
    \item \textbf{Lập kế hoạch}: Dùng AI để breakout các WIGs (Mục tiêu cực kỳ quan trọng) thành các Tasks nhỏ và sắp xếp vào Gantt Chart tự động.
    \item \textbf{Giám sát}: Hệ thống tự động phân tích tiến độ và đưa ra cảnh báo nếu dự án có dấu hiệu chậm trễ.
    \item \textbf{Ghi chép}: Tự động tóm tắt biên bản họp (Minutes of Meeting) và trích xuất các hành động (Action items) cần thực hiện.
\end{itemize}

\subsection{AI trong Phát triển bản thân (Tư duy của một ``Trainee'')}
\begin{itemize}
    \item \textbf{Học tập chủ động}: Thay vì chờ Mentor giải đáp, Mentee được hướng dẫn cách dùng AI làm ``Gia sư 24/7'' để tự đào sâu kiến thức.
    \item \textbf{Tối ưu hóa hiệu suất}: Sử dụng AI để lập lịch làm việc, quản lý mức năng lượng và phản tư (reflection) sâu sắc hơn qua mỗi ngày.
\end{itemize}

\section{Chiến lược Mentoring dành cho Mentor trẻ}

Làm thế nào để một Mentor trẻ ``có tuổi đời chưa bằng tuổi nghề'' của người khác có thể dẫn dắt hiệu quả?

\begin{enumerate}
    \item \textbf{Mentoring bằng thực thi (Show, don't just tell)}: Thay vì dùng triết lý, hãy dùng kết quả. Tôi cùng Mentee xây dựng sản phẩm thật, chạy thật bằng các công cụ AI mới nhất. Sự hào hứng từ việc ``ra kết quả ngay'' chính là chất keo kết dính mạnh nhất.
    \item \textbf{Lấy tốc độ bù kinh nghiệm}: AI giúp chúng ta thực hiện các thử nghiệm nhanh hơn gấp 10 lần. Mentor trẻ cần dạy Mentee cách ``sai nhanh, học nhanh, sửa nhanh''.
    \item \textbf{Xây dựng hệ thống thay vì quản trị con người}: Thay vì liên tục nhắc nhở, tôi xây dựng hệ thống (như WebApp IMP) để AI và Automation làm việc đó. Mentor chỉ xuất hiện ở những ``điểm chạm'' thực sự cần sự thấu cảm và định hướng con người.
\end{enumerate}

\textit{Kết luận}: Mentor trẻ không cần phải ``già đi'' để giỏi hơn, mà cần ``thông minh hơn'' bằng cách đứng trên vai những ``người khổng lồ'' AI.
